\documentclass[10pt]{article}
\usepackage[T1]{fontenc}
\usepackage[utf8]{inputenc}
\usepackage[french]{babel}
\usepackage{amsmath,amssymb,mathrsfs,mathpazo}
\usepackage{geometry}
\geometry{a4paper,margin=2cm}
\usepackage{graphicx,array,tabularx,tabularray,colortbl,multirow,multicol,booktabs,enumitem,pifont,ulem,fancybox,diagbox,cancel}
\usepackage{pstricks,pstricks-add,pst-plot,pst-tree,pst-node,pst-text,pst-func,pst-math,pst-bspline,pst-3dplot,pst-eucl,pst-solides3d}
\let\pscalculate\undefined
\usepackage{pst-bezier}
\usepackage{tikz,pgfplots,tkz-tab}
\usetikzlibrary{arrows,arrows.meta,patterns,calc,positioning,trees,decorations.markings,shapes.geometric,shapes.misc,shapes.arrows,intersections,decorations.pathmorphing,decorations.pathreplacing,backgrounds,shadows,babel}
\pgfplotsset{compat=1.18}
\usepgfplotslibrary{fillbetween}
\usepackage[np]{numprint}
\usepackage{listings,makecell,pas-tableur,twoopt}
\let\barre\undefined
\usepackage[pstricks]{tablvar}
\definecolor{darkgreen}{RGB}{0,100,0}
\providecommand{\red}{\color{red}}
\providecommand{\blue}{\color{blue}}
\providecommand{\green}{\color{green}}
\providecommand{\black}{\color{black}}
\providecommand{\white}{\color{white}}
\providecommand{\gray}{\color{gray}}
\providecommand{\cyan}{\color{cyan}}
\providecommand{\magenta}{\color{magenta}}
\providecommand{\yellow}{\color{yellow}}
\providecommand{\violet}{\color{violet}}
\providecommand{\orange}{\color{orange}}
\providecommand{\euro}{\texteuro}
\providecommand{\EUR}{\texteuro}
\providecommand{\ieme}{\textsuperscript{e}}
\providecommand{\degres}{\ensuremath{^\circ}}
\providecommand{\degree}{\ensuremath{^\circ}}
\providecommand{\vv}[1]{\overrightarrow{#1}}
\providecommand{\Coord}[1]{#1}
\providecommand{\decosix}{}
\providecommand{\starredbullet}{\ensuremath{\bullet}}
\pagestyle{empty}
\setlength{\parindent}{0pt}
\usepackage{adjustbox}\DeclareUnicodeCharacter{2236}{:}\DeclareUnicodeCharacter{2190}{\ensuremath{\leftarrow}}
\providecommand{\R}{\mathbb{R}}
\providecommand{\N}{\mathbb{N}}
\providecommand{\Z}{\mathbb{Z}}
\providecommand{\Q}{\mathbb{Q}}
\providecommand{\D}{\mathbb{D}}
\providecommand{\C}{\mathbb{C}}
\providecommand{\e}{{\mathrm{e}}}
\providecommand{\vect}[1]{\overrightarrow{#1}}
\providecommand{\vectt}[1]{\overrightarrow{\mathrm{#1}}}
\providecommand{\vvt}[1]{\overrightarrow{\mathrm{#1}}}
\providecommand{\np}[2][]{{#2~#1}}
\providecommand{\vv}[1]{{\overrightarrow{#1}}}
\renewcommand{\d}{{\mathrm{d}}}
\providecommand{\degres}{{^\circ}}
\providecommand{\degre}{{^\circ}}
\providecommand{\blue}{\color{blue}}
\providecommand{\red}{\color{red}}
\providecommand{\enskip}{\quad}
\providecommand{\llbracket}{\lbrack\!\lbrack}
\providecommand{\rrbracket}{\rbrack\!\rbrack}
\providecommand{\euro}{{\text{€}}}
\providecommand{\no}{{\text{n°}}}
\providecommand{\interval}[2]{{[#1;#2]}}
\providecommand{\intervFO}[2]{{[#1;#2[}}
\providecommand{\dotfill}{\ldots}

\providecommand{\ds}{\displaystyle}
\providecommand{\textding}[1]{#1}
\providecommand{\Oij}{\(\mathrm{O};\vec{\imath},\vec{\jmath}\)}
\providecommand{\Oijk}{\(\mathrm{O};\vec{\imath},\vec{\jmath},\vec{k}\)}
\providecommand{\Ouv}{\(\mathrm{O};\vec{u},\vec{v}\)}
\providecommand{\intervFF}[2]{\ensuremath{[#1;#2]}}
\renewcommand{\euro}{\ifmmode\text{\texteuro}\else\texteuro\fi}
\DeclareUnicodeCharacter{00B0}{\ifmmode^{\circ}\else\textdegree\fi}
\DeclareUnicodeCharacter{20AC}{\ifmmode\text{\texteuro}\else\texteuro\fi}
% Énoncés : Éducation nationale. Transcription et corrigés : APMEP, Denis Vergès et contributeurs. Découpage et classement : CoopMaths, Sébastien Lozano, Éric Elter et contributeurs. Adaptation de la présentation : Maths Explorer.
% CC BY-SA : https://coopmaths.fr/www/about/
% https://coopmaths.fr/alea/static/bac/2025/tex/bac_2025_09_sujet1_metropole_2.tex
\begin{document}
\noindent\textit{APMEP / CoopMaths. CC BY-SA. Présentation adaptée par Maths Explorer.}\par\medskip
\medskip

\begin{minipage}{0.43\linewidth}
On considère le cube ABCDEFGH.\\

On place le point M tel que $\vectt{BM} = \vectt{AB}$.
\end{minipage}
\hfill
\begin{minipage}{0.55\linewidth}
\psset{unit=1cm}
\begin{pspicture}(7,5.5)
\pspolygon(0.3,0.8)(3.3,0.5)(3.3,3.5)(0.3,3.8)
\psline(3.3,0.5)(6.3,0.2)(3.3,3.5)(5.5,4.7)(2.2,5)(0.3,3.8)
\psline[linestyle=dashed](3.3,0.5)(5.5,1.7)(2.2,2)(0.3,0.8)
\psline[linestyle=dashed](2.2,2)(2.2,5)
\psline(5.5,1.7)(5.5,4.7)
\psline(5.13,1.5)(5.5,1.7)(4.9,1.77)
\uput[dl](0.3,0.8){A} \uput[d](3.3,0.5){B} \uput[dr](5.5,1.7){C}
\uput[ur](2.2,2){D} \uput[l](0.3,3.8){E} \uput[u](3.3,3.5){F}
\uput[ur](5.5,4.7){G} \uput[u](2.2,5){H} \uput[r](6.3,0.2){M}
\end{pspicture}
\end{minipage}

\medskip

\textbf{Partie A}

\medskip

\begin{enumerate}
\item Montrer que les droites (FG) et (FM) sont perpendiculaires.
\item Montrer que les points A, M, G et H sont coplanaires.
\end{enumerate}

\medskip

\textbf{Partie B}

\medskip

On se place dans le repère orthonormé $\left(\text{A} ; \vect{\text{AB}}, \vect{\text{AD}}, \vect{\text{AE}}\right)$.


\medskip

\begin{enumerate}
\item Déterminer les coordonnées des vecteurs $\vect{\text{GM}}$ et $\vect{\text{AH}}$ et montrer qu'ils ne sont pas colinéaires.
\item
	\begin{enumerate}
		\item Justifier qu'une représentation paramétrique de la droite (GM) est :

\[\left\{\begin{array}{l c l}
x&=&1 + t\\
y&=&1 - t\\
z&=&1 - t
\end{array}\right. \:\text{ avec}\: t \in \R.\]

		\item On admet qu'une représentation paramétrique de la droite (AH) est :

\[\left\{\begin{array}{l c l}
x&=&0\\
y&=&k\\
z&=&k
\end{array}\right. \:\text{ avec}\: k \in \R.\]

Montrer que le point d'intersection de (GM) et (AH), que l'on nommera N, a pour
coordonnées (0~;~2~;~2).
	\end{enumerate}
\item
	\begin{enumerate}
		\item Montrer que le triangle AMN est un triangle rectangle en A.
		\item Calculer l'aire de ce triangle.
	\end{enumerate}
\item Soit J le centre de la face BCGF{}.
	\begin{enumerate}
		\item Déterminer les coordonnées du point J.
		\item Montrer que le vecteur $\vect{\text{FJ}}$ est un vecteur normal au plan (AMN).
		\item Montrer que J appartient au plan (AMN). En déduire qu'il est le projeté orthogonal du point F sur le plan (AMN).
	\end{enumerate}
\item On rappelle que le volume $\mathcal{V}$ d'un tétraèdre ou d'une pyramide est donné par la formule :
\[\mathcal{V} = \dfrac13 \times \mathcal{B} \times h,\]

$\mathcal{B}$ étant l'aire d'une base et $h$ la hauteur relative à cette base.

Montrer que le volume du tétraèdre AMNF est le double du volume de la pyramide BCGFM.
\end{enumerate}

\bigskip

\end{document}
